\chapter{BPCC：身体预测控制小脑——SGC认知孔径、世界模型形成与高层预测约束}

\begin{chapteroverviewbox}
本章在前述 Body Runtime、SGC、FCS 与 BPCC 基础上进一步讨论一个容易被忽略、但对于 AgentOS 高层认知形成极其重要的问题：Body Runtime 所维护的现实状态是否应当以完整、高分辨率的形式持续暴露给 Agent Kernel。直观上，向认知内核提供越完整的世界信息似乎越有利于智能，但从工作记忆有限容量、世界模型形成、预测学习以及大小脑职责分离来看，这一设计实际上可能产生相反结果。若高层认知在任何时刻都能够无成本读取完整、高清、持续更新的 SGC，那么系统很容易退化成一个依赖外部输入不断重新计算的反应系统，而不是主动维持世界状态、预测不可见现实并形成高阶结构模型的持续智能体。

因此，本章提出一个新的 Body Runtime 原则：Body Runtime 可以尽可能完整地维护当前身体可获得的现实语义状态，但 Agent Kernel 默认不应持续获得这一完整状态，而只接收经过目标、关注、风险、动态变化和认知容量共同调节的局部高分辨率 SGC 投影。该机制可以类比生物视觉中的中央凹结构：注意中心具有较高几何与语义分辨率，外围区域保持粗粒度对象、运动、风险与变化信息，而异常事件拥有突破当前注意孔径并请求认知升级的能力。

这种设计首先能够降低工作记忆 $h(t)$ 的关系负载，但这并不是它最重要的意义。更深层的作用是人为建立一个适度的认识论瓶颈，使 Agent Kernel 无法简单依赖当前输入获得完整世界，而必须依赖 $h$、$E$ 与 $\Theta$ 共同维持一个跨时间连续的内部世界模型。由此，高层认知被迫学习对象恒常性、空间连续性、事件发展、隐藏状态、因果关系和未来预测，而不是把有限的工作记忆容量浪费在高频、低层、短寿命的几何与控制细节上。

与此同时，这种输入约束还直接决定了 Kernel 向 BPCC 输出何种信息。若 Kernel 接收大量低层控制细节，它便容易自然地产生低层微操作指令，使 BPCC 退化成简单执行器。相反，如果 Kernel 的主要输入本身已经经过结构抽象和有限观测约束，那么它更容易形成对象级、关系级、任务级和未来状态级的预测，并将这些预测压缩成 Predictive Envelope、Control Reference 与 Semantic Constraint 发送给 BPCC。于是大小脑之间形成真正的多尺度分工：Kernel 负责预测世界将向什么高层状态演化，BPCC 负责利用自身完整、高频的局部现实状态决定身体具体怎样到达这些状态。

本章因此把 SGC 输入控制、$h(t)$ 世界模型、主动感知、高层预测与 BPCC 快速控制统一进同一个多尺度闭环，并提出一个核心原则：

\begin{statementbox}
身体持续维护足够广泛的现实，高层认知只被允许看清有限的一部分，
\end{statementbox}

\begin{statementbox}
从而迫使认知理解未被直接看见的世界，并将这种理解转换成对快速身体闭环的未来结构约束。
\end{statementbox}
\end{chapteroverviewbox}

\section{完整现实状态与认知现实输入必须分离}

在早期设计中，可以简单写成：

\[
World
\rightarrow
SGC
\rightarrow
AgentKernel.
\]

但随着 SGC 内容不断丰富，这一表达已经不够准确。SGC 不仅包含视觉对象，还可能包括三维几何、实体身份、固定属性、语义标签、对象关系、场信息、身体位置、Affordance、Epistemic Metadata 以及当前预测结果。如果 Body Runtime 将这些内容全部以最高分辨率持续输入 Kernel，那么高层认知面对的实际上不是一个有限的认知窗口，而是一个近似无边界的高带宽世界数据库。

因此应严格区分三个层次。第一层是 Body Runtime 当前维护的尽可能完整的身体相关现实状态：

\[
\boxed{
\mathcal G_t^{full}=SGC_{full}(t)
}
\]

其中包含 Body Runtime 当前有能力维护的实体、场、几何、关系和状态。这里所谓 ``full'' 并不意味着宇宙真实状态的完整副本，而只是相对于当前身体感知和局部世界模型而言的最大可用 SGC。

第二层是认知观察状态：

\[
\boxed{
A_t=
(c_t,\Omega_t,s_t,q_t,\chi_t)
}
\]

其中 $c_t$ 表示当前关注中心，$\Omega_t$ 表示关注区域，$s_t$ 表示观察尺度，$q_t$ 表示分辨率分布，$\chi_t$ 表示任务或语义筛选条件。

第三层才是真正进入 Agent Kernel 的 SGC：

\[
\boxed{
SGC_{view}(t)
=
\mathcal R
\left(
SGC_{full}(t),A_t
\right)
}
\]

其中 $\mathcal R$ 是 SGC Cognitive Projection 或 SGC Foveated Rendering 算子。

因此必须建立：

\[
\boxed{
SGC_{full}\neq SGC_{view}
}
\]

这一工程原则。Body Runtime 可以持续拥有较广泛的现实，但认知内核不应永久拥有同等分辨率的整个现实。

这种区分实际上把 Body Runtime 从单纯的感知编码器提升成了一个主动认知前端。它不仅回答 ``现实里有什么''，还需要回答：

\begin{statementbox}
现实当前应当以什么结构分辨率进入认知？
\end{statementbox}

这正是有限认知系统必须解决的第一道问题。

\section{SGC中央凹不是图像模糊，而是结构分辨率变化}

如果 SGC 只是 RGB 图像，那么 ``中央高清、边缘模糊'' 可以通过传统视觉分辨率直接理解。但 SGC 是结构化的世界表示，因此其中央凹机制必须发生在语义结构层，而不能仅仅作用于像素。

设当前注意中心为 $c_t$，可以定义空间分辨率函数：

\[
q(x\mid c_t)
=
q_{\max}
e^{-\alpha d(x,c_t)}
+
q_{\min},
\]

其中 $d(x,c_t)$ 表示某个空间位置或实体相对于当前关注中心的距离。随着距离增大，空间与结构表示精度逐渐下降：

\[
d(x,c_t)\uparrow
\Rightarrow
q(x\mid c_t)\downarrow.
\]

但这种下降不仅意味着几何坐标精度降低，也意味着实体属性、语义关系和预测细节逐渐被压缩。例如一个中央区域中的杯子可以被表示为：

\[
\{
Identity,
FineGeometry,
Material,
Pose,
FillLevel,
Relations,
Affordance,
Prediction
\},
\]

而在外围只保留：

\[
\{
ObjectClass,
CoarsePosition,
Motion,
Risk
\},
\]

更远处甚至可以压缩为：

\[
\{
Occupancy,
ObjectCluster,
Saliency
\}.
\]

因此可以建立结构分辨率层级：

\[
L_0: Occupancy,
\]

\[
L_1: EntityClass,
\]

\[
L_2: Entity+Attributes,
\]

\[
L_3:
Entity+Attributes+Relations+FineGeometry+Semantics.
\]

高层认知真正受到控制的是：

\[
\boxed{
ResolutionOfAgentCSOStructure
}
\]

而不只是感知像素数量。

这种结构分辨率机制使 SGC 比生物视觉中央凹更加一般，因为它可以同时控制空间、时间和语义三个方向的分辨率，而不是只有二维视网膜上的空间清晰度变化。

\section{三维认知分辨率：空间、语义与时间}

SGC 的认知孔径至少应同时作用于三个维度。

第一是空间分辨率：

\[
R_x(o_i),
\]

决定对象 $o_i$ 的几何坐标、形状、姿态和局部环境以多大精度进入 Kernel。

第二是语义分辨率：

\[
R_s(o_i),
\]

决定对象保留多少属性、关系和高层语义。一个离当前视线很远但与当前 Goal 强相关的对象，仍然可以保持高语义分辨率，因此语义分辨率不能简单由空间距离决定。

第三是时间分辨率：

\[
R_t(o_i),
\]

决定对象状态以多高频率更新。高速运动物体、快速变化风险或短时间内可能影响当前任务的实体，即使位于认知外围，也可能需要保持较高的时间更新频率。

因此一个更加一般的分辨率控制律可以写成：

\[
\boxed{
R_i
=
f(
D_i,
G_i,
V_i,
Risk_i,
U_i,
C_h
)
}
\]

其中 $D_i$ 表示空间距离，$G_i$ 表示目标相关度，$V_i$ 表示动态变化速度，$Risk_i$ 表示风险，$U_i$ 表示预测不确定性，而 $C_h$ 表示当前工作记忆可用容量。

由此，SGC 中央凹不是固定圆形区域，而是一个根据任务与动力学不断变化的：

\[
\boxed{
Semantic-Epistemic Aperture.
}
\]

\section{SGC认识孔径与有益感知稀缺原则}

这里可以正式定义 SGC Epistemic Aperture，简称 SEA：

\[
\boxed{
SEA=
SGC\ Epistemic\ Aperture.
}
\]

SEA 接收：

\[
SGC_{full},
\]

以及：

\[
AttentionReference,
Goal,
Risk,
Dynamics,
Capacity,
\]

输出：

\[
SGC_{view}.
\]

形式上：

\[
\boxed{
SGC_{view}
=
SEA
\left(
SGC_{full},
Goal,
Attention,
Risk,
Dynamics,
C_h
\right)
}
\]

并要求：

\[
InformationDensity(SGC_{view})
\ll
InformationDensity(SGC_{full}),
\]

但同时尽量保持：

\[
TaskRelevantInformation(SGC_{view})
\approx
TaskRelevantInformation(SGC_{full}).
\]

这一机制不应仅被理解为信息压缩，因为它实际上有意识地向 Kernel 保留一个认识论缺口。因此可以进一步提出：

\[
\boxed{
Principle\ of\ Productive\ Perceptual\ Scarcity
}
\]

即：

\begin{statementbox}
有益感知稀缺原则
\end{statementbox}

该原则认为，对于高层智能而言，最优输入并不一定是最大信息输入。高层认知应获得足以支持当前任务，却不足以让其完全依赖即时输入的现实表示。这样才能迫使内部状态承担记忆、预测和结构抽象职责。

可以概念性写成：

\[
\min_{I_K}
\left[
CognitiveLoad(I_K)
+
\lambda D_{input}(I_K)
\right]
\]

约束：

\[
Performance(I_K)\ge P_{\min},
\]

其中 $D_{input}$ 表示 Kernel 对即时外部输入的依赖程度。

因此：

\[
\boxed{
MoreObservation
\not\Rightarrow
BetterHighLevelIntelligence.
}
\]

对高层认知而言，适度缺失反而可以成为促进内部建模的重要结构条件。

\section{局部高清输入强迫h形成真正的世界模型}

如果 Kernel 每个周期都可以读取完整 SGC，那么最容易形成的行为模式是：

\[
Observation_t
\rightarrow
Inference_t
\rightarrow
Action_t.
\]

系统可以不断依靠重新观察现实解决当前问题，而没有足够压力去维护持续世界状态。此时 $h(t)$ 很容易退化成一个实时输入缓存和临时推理工作区。

引入局部高清 SGC 后：

\[
SGC_{view,t}
\subset
SGC_{full,t},
\]

高层若要保持任务连续，就必须构造：

\[
\boxed{
\hat W_t
=
F_W
\left(
SGC_{view,t},
h_{t-1},
E_t,
\Theta_t
\right)
}
\]

其中 $\hat W_t$ 表示 Agent 当前内部维持的世界估计。

此时系统开始出现一个非常重要的区分：

\[
\boxed{
ObservedWorld
\neq
BelievedWorld
}
\]

其中 ObservedWorld 是当前孔径内直接获得的信息，而 BelievedWorld 是 Agent 根据当前观测、历史经验和生成结构共同维持的内部世界。

真正持续智能必须能够做到：

\[
\boxed{
WorldContinuesInsideAgentWhenNotObserved.
}
\]

例如，一个杯子离开当前中央高清区域以后，并不意味着它从 Agent 世界中消失。系统应保持：

\[
P(
CupPosition_{t+\Delta}
\mid
CupPosition_t,
WorldDynamics,
Context
).
\]

因此：

\[
\boxed{
Visible\neq Exists.
}
\]

对象恒常性、空间连续性和隐藏状态预测由此成为 Kernel 必须主动学习的能力，而不再能够依赖传感器每个周期重新回答。

\section{有限观测使预测误差成为真正的学习信号}

局部高清机制还有一个更深的作用：它人为创造了预测必须承担责任的区域。

在当前注意区域之外，Kernel 不能直接读取完整真实状态，因此只能产生：

\[
\hat{SGC}_{t+\Delta}^{K}.
\]

当稍后重新观察该区域时，真实观测可以与此前预测比较：

\[
\boxed{
\varepsilon_t^{K}
=
SGC_{observed,t}
-
\hat{SGC}_{predicted,t}.
}
\]

这里的减法是结构意义上的差异，并不要求所有 SGC 内容都具有线性数值形式。误差可以包括：

位置误差；

实体身份误差；

关系误差；

事件发展误差；

属性变化误差；

对象出现/消失误差；

风险预测误差。

因此高层形成与 BPCC 类似的基本预测循环：

\[
\boxed{
Prediction
\rightarrow
Observation
\rightarrow
Residual
\rightarrow
Update.
}
\]

不同的是，BPCC 处理的是高速局部物理状态，而 Kernel 处理的是更慢、更广泛的语义与任务状态。

可以写成：

\[
BPCC:
\quad
\hat x_{t+\delta}
\rightarrow
x_{t+\delta}
\rightarrow
\varepsilon_{local},
\]

而：

\[
Kernel:
\quad
\hat{\mathcal G}_{t+\Delta}
\rightarrow
\mathcal G_{t+\Delta}
\rightarrow
\varepsilon_{semantic}.
\]

其中：

\[
\delta\ll\Delta.
\]

由此形成真正的跨时间尺度预测系统。

\section{有限感知迫使h寻找跨时间稳定结构}

如果所有低层细节都能够长期稳定地进入 $h$，那么 Kernel 没有充分动力进行抽象。它可以直接利用大量精确坐标、瞬时属性和当前传感器状态完成推理。

但当 SGC 的局部高清区域不断移动时，高频细节变得暂时和不可靠，而只有某些更稳定的变量能够跨多个注意窗口维持：

\[
ObjectIdentity,
\]

\[
SpatialRelation,
\]

\[
CausalRelation,
\]

\[
GoalState,
\]

\[
EventStructure,
\]

\[
Affordance,
\]

\[
InvariantProperty.
\]

因此局部高清输入形成了一种：

\[
\boxed{
StructuralAbstractionPressure.
}
\]

系统如果希望保持跨时间一致性，就必须逐渐从：

\[
ExactState
\]

转向：

\[
InvariantStructure.
\]

例如 BPCC 可以持续维护杯子的精确姿态：

\[
p_t,\theta_t,\dot p_t,
\]

而 $h$ 更适合保留：

\[
\boxed{
Cup\ is\ stably\ located\ on\ the\ target\ table.
}
\]

这种结构才适合继续进入：

\[
E
\]

并最终沉降到：

\[
\Theta.
\]

因此 SEA 不只是减少认知负担，也成为推动三相记忆世界模型形成的重要训练条件。

\section{SGC局部高清与三相记忆世界模型}

局部观测约束能够直接推动：

\[
h
\rightarrow
E
\rightarrow
\Theta
\]

的结构形成。

在 $h$ 中，Agent 维护当前世界假设：

\[
h_t:
\quad
CurrentEstimatedWorld.
\]

当这些预测在多次观测中得到验证，它们进入 E：

\[
E:
\quad
ExperiencedWorldTrajectories.
\]

重复经验进一步形成 $\Theta$ 中稳定的世界生成结构，例如：

\[
\text{物体离开当前视野后通常仍持续存在},
\]

\[
\text{运动物体具有近似连续的轨迹},
\]

\[
\text{墙体遮挡不意味着墙后空间不存在},
\]

\[
\text{人的行动具有短时间目标连续性}.
\]

于是：

\[
\boxed{
PartialObservation
\rightarrow
Prediction
\rightarrow
Experience
\rightarrow
GenerativeStructure.
}
\]

因此 SGC Foveation 可以被理解为：

\[
\boxed{
WorldModelTrainingPressure.
}
\]

它通过不给 Kernel 完整答案，迫使 Kernel 形成能够补充现实缺口的 $\Theta$。

\section{Kernel认知注意与BPCC控制注意必须严格分离}

本章最重要的工程边界之一，是局部高清机制只能约束 Kernel 的认知输入，而不能成为 BPCC 的控制输入限制。

必须满足：

\[
\boxed{
BPCCInput
\neq
SGC_{view}.
}
\]

BPCC 应使用自己需要的高速控制状态：

\[
\boxed{
I_{BPCC}
=
SGC_{control}
\oplus
FCS
\oplus
z_{BPCC}
\oplus
ControlReference.
}
\]

其中：

\[
SGC_{control}
\]

表示局部控制所需要的高频现实结构，

\[
z_{BPCC}
\]

表示 BPCC 的私有 latent 状态。

Kernel 的输入则是：

\[
\boxed{
I_K
=
SGC_{view}
\oplus
E
\oplus
\Theta
\oplus
\Psi
\oplus
Goal.
}
\]

因此应该严格建立：

\[
\boxed{
CognitiveAttention
\neq
ControlAttention.
}
\]

例如机器人正在端一杯水并与用户交谈。Kernel 的认知焦点可能位于：

\[
UserFace,
Conversation,
Goal,
CupDelivery,
\]

而 BPCC 的控制焦点则位于：

\[
JointState,
GripperForce,
CupPose,
CenterOfMass,
NearbyObstacle.
\]

高层当前没有 ``关注'' 某个关节，并不意味着身体控制器可以停止监控该关节。

因此必须遵循：

\[
\boxed{
FoveationAffectsCognition,
NotRealtimeControlClosure.
}
\]

\section{FCS不能简单服从SGC认知孔径}

SGC 描述：

\[
\boxed{
WhatIsThere?
}
\]

FCS 描述：

\[
\boxed{
HowIsRealityAffectingMe?
}
\]

二者功能不同，因此不能简单把 SGC 的中央凹机制复制到 FCS。

例如 Agent 当前认知焦点在用户面部，但机器人腿部发生：

ThermalPressure；

IntegrityRisk；

MotorOverload；

CollisionRisk。

这些 FCS 信号必须继续进入：

\[
BPCC,
\]

\[
AHC,
\]

甚至：

\[
SafetyChannel,
\]

而不能因为 Kernel 当前没有注意该身体区域而被降低到不可用程度。

因此：

\[
\boxed{
SGCFoveation
\nRightarrow
FCSFoveation.
}
\]

尤其高风险、高强度 FCS 应具有：

\[
\boxed{
AttentionBypass.
}
\]

这再次说明 Body Runtime 的目标不是模拟一个被动视觉系统，而是维持一个多通道、不同时间尺度和不同权限等级的现实耦合界面。

\section{高层预测不应输出低层动作，而应输出Predictive Envelope}

有限 SGC 输入最重要的控制意义，在于它迫使 Kernel 形成适合传递给 BPCC 的高层未来结构。

如果 Kernel 拥有完整、高频、低层控制信息，它很容易开始输出：

\[
LowLevelState
\rightarrow
LowLevelCommand.
\]

例如直接规定：

\[
Joint_1(t+\delta),
\quad
Joint_2(t+\delta),
\quad
GripForce(t+\delta).
\]

这会使 BPCC 退化成执行器，并破坏大小脑之间的时间尺度分工。

正确关系应当是：

\[
\boxed{
Kernel
\rightarrow
PredictiveEnvelope
\rightarrow
BPCC.
}
\]

Kernel 预测：

\[
\hat{\mathcal G}_{t:t+H}^{K},
\]

但不必把完整未来世界传给 BPCC，而应抽取高层约束：

\[
\boxed{
\mathcal E_K
=
(
GoalRegion,
DesiredRelation,
ForbiddenRegion,
TemporalConstraint,
RiskConstraint,
Priority,
Confidence
).
}
\]

BPCC 则在该约束面内求自己的局部最优控制：

\[
u_{t:t+H}^{*}
=
\arg\min_{u}
J_{BPCC}
\]

subject to

\[
x_{t:t+H}
\in
\mathcal E_K.
\]

因此高层真正控制的不是：

\[
\text{动作轨迹本身},
\]

而是：

\begin{statementbox}
未来允许状态空间的形状。
\end{statementbox}

这就是此前多尺度理论中：

\[
\boxed{
SlowStructureShapesFastDynamics
}
\]

在 Body Runtime 中最直接的工程实现。

\section{反微操原则}

由此可以正式提出 AgentOS 的：

\[
\boxed{
Anti\text{-}Micromanagement\ Principle
}
\]

即反微操原则：

\begin{quote}
Agent Kernel 不应获得足以诱导其长期接管快速身体闭环的大量低层高频状态，也不应主要向 BPCC 发送动作级微操作指令。Kernel 应优先处理跨时间稳定的对象、关系、目标、风险和未来状态结构，并把这些结构压缩成慢尺度预测约束，由 BPCC 在完整局部现实条件下完成高速控制。
\end{quote}

因此：

\[
\boxed{
Resolution(KernelInput)
<
Resolution(BPCCControlState)
}
\]

而在抽象等级上：

\[
\boxed{
Abstraction(KernelPrediction)
>
Abstraction(BPCCPrediction).
}
\]

这样才能真正形成：

\[
\boxed{
Kernel=Meaning/FutureConstraint,
}
\]

\[
\boxed{
BPCC=RealtimePhysicalRealization.
}
\]

\section{从SGC输入限制到BPCC控制的完整因果链}

上述机制使看似独立的 SGC 感知设计和 BPCC 控制设计形成统一链条：

\[
\boxed{
SGCInputRestriction
}
\]

首先减少高层即时现实暴露；

然后：

\[
\boxed{
ForcedInternalWorldModeling
}
\]

迫使 $h/E/\Theta$ 维持当前不可见状态；

进一步：

\[
\boxed{
HighLevelStructuralPrediction
}
\]

形成对象、关系、任务和未来事件级预测；

随后压缩为：

\[
\boxed{
PredictiveEnvelope
}
\]

再传给：

\[
\boxed{
BPCCFastControl.
}
\]

即：

\[
\boxed{
SGC_{full}
\rightarrow
SGC_{view}
\rightarrow
h
\rightarrow
\hat W
\rightarrow
\hat W_{future}
\rightarrow
PredictiveEnvelope
\rightarrow
BPCC
}
\]

而 BPCC 行动重新作用现实：

\[
BPCC
\rightarrow
Action
\rightarrow
World'
\rightarrow
SGC_{full}'.
\]

整个闭环为：

\[
\boxed{
World
\rightarrow
SGC_{full}
\rightarrow
SGC_{view}
\rightarrow
Kernel
\rightarrow
PredictiveEnvelope
\rightarrow
BPCC
\rightarrow
World'.
}
\]

这表明 SGC Foveation 不是感知系统的附属功能，而是 Kernel—BPCC 多尺度闭环中决定信息抽象梯度的重要机制。

\section{中央高清、外围摘要与异常哨兵}

有限 SGC 输入并不意味着中央高清区之外的世界完全消失。正确结构应至少包括三层。

第一层是中央认知区：

\[
\boxed{
FovealRegion
}
\]

保留完整或高分辨率：

Geometry；

Attributes；

Relations；

Semantics；

Prediction；

Affordance。

第二层是外围摘要区：

\[
\boxed{
PeripheralSummary
}
\]

只保留：

粗位置；

对象类别；

运动状态；

风险；

变化事件；

当前置信度。

第三层是异常哨兵：

\[
\boxed{
SentinelChannel.
}
\]

它持续监测：

新实体；

突然运动；

高速接近；

高风险；

高 Goal 相关度；

异常声源；

高预测残差。

因此 Kernel 输入不是：

\[
OnlyFovea,
\]

而是：

\[
\boxed{
KernelInput
=
FovealDetail
+
PeripheralSummary
+
SaliencyInterrupt.
}
\]

这样既保持认知有限性，又避免产生严重注意盲区。

\section{外围结构的注意抢占机制}

设外围对象 $o_i$ 的显著性：

\[
S_i
=
w_rR_i
+
w_nN_i
+
w_dD_i
+
w_gG_i
+
w_uU_i,
\]

其中：

$R_i$ 为 Risk，

$N_i$ 为 Novelty，

$D_i$ 为 Dynamic Change，

$G_i$ 为 Goal Relevance，

$U_i$ 为 Predictive Uncertainty。

如果：

\[
S_i>\theta_{focus},
\]

则该对象进入：

\[
\boxed{
FovealCandidate.
}
\]

随后 Scheduler 或 Kernel 决定是否真正改变：

\[
AttentionReference.
\]

由此形成：

\[
PeripheralDetection
\rightarrow
Saliency
\rightarrow
AttentionRequest
\rightarrow
FocusShift.
\]

这使 SEA 与：

SPC；

TCE；

CCM；

Scheduler

形成新的认知控制闭环。

特别是 CCM 可以根据当前 $h$ 的剩余容量动态调整 SGC 分辨率。如果：

\[
C_h\rightarrow C_{critical},
\]

则可以主动缩窄：

\[
\Omega_t,
\]

降低外围语义分辨率；反之，当认知容量充足且任务需要搜索时，可以扩大认识孔径。

因此：

\[
\boxed{
PerceptualAperture
}
\]

本身成为认知资源控制变量。

\section{聚焦、缩放和选区属于认识动作而不是世界动作}

此前 Body Runtime 中曾讨论过 Focus、Select、Zoom 等只读操作。经过本章重新解释后，这些操作不应主要按照 UI 行为理解，而应归入：

\[
\boxed{
EpistemicAction.
}
\]

普通行动改变：

\[
World.
\]

而认识行动改变：

\[
ObservationOperator.
\]

设：

\[
O_{A_t}(World_t)
\]

为当前观测算子，则：

\[
Focus,
Zoom,
Select
\]

改变的是：

\[
A_t,
\]

从而改变：

\[
O_{A_t},
\]

而不是直接改变：

\[
World_t.
\]

所以：

\[
\boxed{
ReadOnlyControl
=
ControlOfObservation,
NotControlOfWorld.
}
\]

由此以后 AgentOS 的 Action 可以至少区分成三类：

\[
\boxed{
EpistemicAction
}
\]

改变 Agent 看什么、怎样看；

\[
\boxed{
WorldAction
}
\]

直接改变物理或数字外部世界；

\[
\boxed{
HomeostaticAction
}
\]

改变自身身体或内部资源状态。

这种分类对于后续 AgentOS Action ABI、Body ABI 和权限模型都具有直接意义。

\section{主动感知：从“看哪里”发展成认知实验}

一旦观察具有成本和分辨率限制，下一步 ``看哪里'' 本身就成为一个决策问题。

Agent 可以根据当前世界模型选择下一次 Focus：

\[
A_t^*
=
\arg\max_{A}
\mathbb E
\left[
IG(A)
-
\lambda C(A)
\right],
\]

其中 $IG(A)$ 表示预期 Information Gain，$C(A)$ 表示新的观察所占用的认知资源、动作成本或时间成本。

于是 Agent 的感知闭环从：

\[
LookEverywhere
\rightarrow
React
\]

转变成：

\[
\boxed{
Predict
\rightarrow
ChooseWhereToLook
\rightarrow
Observe
\rightarrow
Compare
\rightarrow
Learn.
}
\]

例如 Agent 认为杯子仍然位于桌面：

\[
P(CupOnTable)=0.92.
\]

当某个声音或外围运动使不确定性提高后，它主动把 SEA 中央区域移到桌面进行验证。如果结果与预测一致，则增强相应世界结构；如果不一致，则产生新的 semantic residual，并更新 $h/E/\Theta$。

因此注意不再只是资源分配机制，还成为一种：

\[
\boxed{
ActiveEpistemicExperiment.
}
\]

\section{局部高清输入防止外部世界退化为外挂工作记忆}

如果现实界面永远向 Kernel 提供完整高分辨率信息，那么外部世界本身就可能退化成一个巨大的 External Working Memory。Agent 无需真正维持内部世界，只需反复扫描环境。

例如数字 Agent 可以每一步重新读取所有窗口、所有文本、所有按钮，而不必记住：

窗口结构；

用户意图；

之前发生过什么；

对象是否移动；

事件是否连续。

这种系统在传感器正常时可能表现良好，但一旦出现：

遮挡；

网络延迟；

窗口隐藏；

摄像头暂时不可用；

对象离开视野；

就容易失去世界连续性。

局部高清机制则强迫：

\[
\boxed{
WorldMustContinueInsideTheAgent.
}
\]

因此 Agent 才能逐步形成真正的：

\[
PersistentWorldModel.
\]

这个作用远比单纯节约上下文更加重要。

\section{信息分辨率梯度与AgentOS双向结构流}

结合本章机制，可以观察到整个 AgentOS 中存在一个非常清晰的信息分辨率梯度。

从现实向高层认知：

\[
Reality
\rightarrow
Sensor
\rightarrow
SGC_{full}
\rightarrow
SGC_{view}
\rightarrow
h
\rightarrow
E
\rightarrow
\Theta.
\]

总体趋势是：

\[
\boxed{
Detail
\rightarrow
Structure,
}
\]

\[
\boxed{
Fast
\rightarrow
Slow,
}
\]

\[
\boxed{
Local
\rightarrow
General,
}
\]

\[
\boxed{
Observation
\rightarrow
Model.
}
\]

反方向则是：

\[
\Theta
\rightarrow
h
\rightarrow
PredictiveEnvelope
\rightarrow
BPCC
\rightarrow
Actuator.
\]

总体趋势是：

\[
\boxed{
Structure
\rightarrow
Constraint,
}
\]

\[
\boxed{
Slow
\rightarrow
Fast,
}
\]

\[
\boxed{
General
\rightarrow
LocalRealization.
}
\]

因此整个 AgentOS 可以进一步抽象成一个双向多尺度结构流：

\[
\boxed{
Upward:
RealityDetail
\rightarrow
CognitiveStructure
}
\]

以及：

\[
\boxed{
Downward:
CognitiveStructure
\rightarrow
FutureConstraint
\rightarrow
PhysicalAction.
}
\]

SGC Foveation 正处在第一条流的关键入口，而 BPCC 位于第二条流的最后快速闭环。

\section{从CSO角度重新解释SGC认知孔径}

如果进一步结合前文 CSO 与 Agent-CSO 理论，SGC 认知孔径还有一个更基础的本体意义。

外部现实中存在大量：

\[
UniverseCSO.
\]

Body Runtime 根据身体当前可感知范围形成：

\[
SGC_{full}.
\]

SEA 决定其中哪些结构以较高分辨率进入 Agent Kernel：

\[
UniverseCSO
\rightarrow
SGC_{full}
\rightarrow
SGC_{view}
\rightarrow
AgentCSO.
\]

因此 SEA 实际控制的是：

\[
\boxed{
AgentCSOResolution.
}
\]

Agent 并不是简单处于：

\[
Seen/Unseen
\]

二值状态，而是可能以不同分辨率承载同一个现实结构：

\[
AgentCSO^{fine},
\quad
AgentCSO^{coarse},
\quad
AgentCSO^{predicted}.
\]

例如外围只知道：

\[
PersonExistsNearDoor,
\]

中央高清后才得到：

\[
Person=
User_A,
\]

\[
Pose=Standing,
\]

\[
Expression=Concerned,
\]

\[
Relation=WaitingForCup.
\]

这说明认知不是现实结构的全部复制，而是：

\[
\boxed{
SelectiveMultiresolutionRealizationOfReality.
}
\]

因此 SGC Foveation 也可以被理解成 Universe-CSO 到 Agent-CSO 之间的一种动态承载映射控制机制。

\section{BPCC为何必须保持独立局部现实闭合}

本章必须再次强调：高层认知的有益感知稀缺不能被误推广为身体控制的信息稀缺。BPCC 的存在本身就是为了避免高层认知承担高频现实闭合。

BPCC 应满足：

\[
\boxed{
NormalBodyDynamicsStayLocal.
}
\]

其局部闭环为：

\[
SGC_{control,t}
+
FCS_t
+
z_t
+
Envelope_t
\rightarrow
Prediction
\rightarrow
Control
\rightarrow
Reality
\rightarrow
SGC_{control,t+\delta}.
\]

只有当局部残差长期不能吸收，或者：

\[
Residual>\theta_{escalation},
\]

才向 Kernel 输出：

\[
Exception,
\quad
ResidualSummary,
\quad
CapabilityViolation.
\]

因此：

\[
\boxed{
KernelAttention
}
\]

不决定身体是否继续看见某个控制变量。

换句话说：

\[
\boxed{
CognitionMayLookAway;
ControlMayNotLoseReality.
}
\]

这是整个 Body Runtime 设计必须严格保证的安全原则。

\section{SGC认知孔径与工作记忆有限容量定理的关系}

此前工作记忆有限容量定理指出：

\[
C_{run}(t)
\le
C_{\max}(W).
\]

CCM 负责进入认知后的结构治理，Boundary 负责决定什么值得进入，而 SEA 则把治理进一步前移到现实表示层。

因此可以建立四级输入防线：

\[
\boxed{
World
}
\]

\[
\downarrow
\]

\[
\boxed{
SGC\ Foveation
}
\]

控制现实的结构分辨率；

\[
\downarrow
\]

\[
\boxed{
Boundary
}
\]

决定哪些结构获得认知准入；

\[
\downarrow
\]

\[
\boxed{
CCM
}
\]

管理已进入系统的结构负载；

\[
\downarrow
\]

\[
\boxed{
h(t)
}
\]

作为有限全局认知工作区。

因此：

\[
\boxed{
SGC\ Foveation
\neq
Boundary
\neq
CCM.
}
\]

三者分别回答：

\begin{statementbox}
现实看多清？
\end{statementbox}

\begin{statementbox}
哪些东西值得进来？
\end{statementbox}

\begin{statementbox}
进来以后怎样保持可处理？
\end{statementbox}

三者共同保护有限 $h(t)$，但其功能不可相互替代。

\section{本章的统一动力学模型}

综合本章，可以把 Kernel 与 BPCC 的现实耦合写成两个并行而交互的闭环。

认知闭环：

\[
\boxed{
SGC_{full}
\rightarrow
SEA
\rightarrow
SGC_{view}
\rightarrow
h/E/\Theta
\rightarrow
WorldPrediction
\rightarrow
AttentionReference
\rightarrow
SEA
}
\]

其中 Kernel 根据当前内部世界模型主动决定下一步需要观察什么。

身体快速闭环：

\[
\boxed{
SGC_{control}
+
FCS
+
PredictiveEnvelope
\rightarrow
BPCC
\rightarrow
Action
\rightarrow
World
\rightarrow
SGC_{control}'.
}
\]

两者之间通过高层预测和异常残差交互：

\[
Kernel
\xrightarrow{
PredictiveEnvelope
}
BPCC,
\]

\[
BPCC
\xrightarrow{
Residual/Exception
}
Kernel.
\]

而 Kernel 的预测来自：

\[
\boxed{
SGC_{view}
+
h
+
E
+
\Theta.
}
\]

因此完整结构为：

\[
\boxed{
\begin{aligned}
World
&\rightarrow
SGC_{full}
\rightarrow
SGC_{view}
\rightarrow
h/E/\Theta\\
&\rightarrow
SemanticPrediction
\rightarrow
PredictiveEnvelope
\rightarrow
BPCC\\
&\rightarrow
Action
\rightarrow
World'.
\end{aligned}
}
\]

同时：

\[
World
\rightarrow
SGC_{control}/FCS
\rightarrow
BPCC
\]

保持独立高频现实闭合。

这意味着整个架构并不存在单一 ``Agent 所看到的世界''，而至少存在三个不同功能层面的现实状态：

\[
\boxed{
SGC_{full}
}
\]

Body Runtime 当前维护的广泛现实；

\[
\boxed{
SGC_{view}
}
\]

Kernel 当前高层认知真正获得的现实；

\[
\boxed{
SGC_{control}
}
\]

BPCC 当前高速局部控制必须知道的现实。

三者来自同一个现实，但服务不同的功能与时间尺度。

\section{本章基本原则}

基于上述分析，可以把本章形成的设计原则正式总结如下。

第一，\textbf{完整现实维护原则}：

\[
\boxed{
BodyRuntime
\text{应尽可能维持广泛的身体相关现实状态。}
}
\]

第二，\textbf{认知输入有限原则}：

\[
\boxed{
Kernel
\text{默认不直接获得完整高分辨率 }SGC_{full}.
}
\]

第三，\textbf{有益感知稀缺原则}：

\begin{statementbox}
高层感知应足以完成任务，但不应丰富到使内部世界模型失去必要性。
\end{statementbox}

第四，\textbf{世界连续性原则}：

\[
\boxed{
NotObserved
\neq
NotExisting.
}
\]

第五，\textbf{预测优先原则}：

\[
\boxed{
Kernel
\text{应通过有限观测形成并维护世界预测，而不是持续依赖实时重读。}
}
\]

第六，\textbf{高层结构优先原则}：

\[
\boxed{
h
\text{主要维持对象、关系、目标、事件和未来结构，而非高频局部控制细节。}
}
\]

第七，\textbf{认知注意—控制注意分离原则}：

\[
\boxed{
Attention_{cognition}
\neq
Attention_{control}.
}
\]

第八，\textbf{FCS独立调制原则}：

\[
\boxed{
SGCFoveation
\nRightarrow
CriticalFCSFoveation.
}
\]

第九，\textbf{反微操原则}：

\[
\boxed{
Kernel
\text{定义未来状态约束，而不是持续规定身体高速轨迹。}
}
\]

第十，\textbf{预测包络原则}：

\[
\boxed{
Kernel
\rightarrow
PredictiveEnvelope
\rightarrow
BPCC.
}
\]

第十一，\textbf{局部现实权威原则}：

\[
\boxed{
BPCC
\text{必须使用自己的实时局部现实状态完成闭环。}
}
\]

第十二，\textbf{异常升级原则}：

\[
\boxed{
LocalResidual
\rightarrow
GlobalCognitiveEscalation
}
\]

只在局部模型不能解决问题时发生。

\section{本章结论：限制看见，才能真正学会理解}

本章最容易被表面理解为一个用于降低上下文和计算量的感知优化方案，但这种理解是不完整的。局部高清 SGC 的首要工程收益当然包括降低 Kernel 输入规模、关系数量和预测空间，但它更重要的作用，是塑造 Agent Kernel 的认知方式。

如果系统永远拥有完整即时现实，高层最容易学习的是：

\[
\boxed{
Look\rightarrow React.
}
\]

而如果现实以有限、高分辨率局部窗口进入高层，同时外围保持粗摘要和异常哨兵，则系统被迫学习：

\[
\boxed{
Observe
\rightarrow
Remember
\rightarrow
Predict
\rightarrow
ChooseWhatToObserve
\rightarrow
Compare
\rightarrow
Update.
}
\]

由此 $h(t)$ 才真正成为世界状态的主动维持面，$E$ 才积累跨时间事件结构，$\Theta$ 才能形成稳定世界生成规律。更重要的是，高层输出也因此自然从低层动作微操上升为：

\[
Goal,
\quad
Relation,
\quad
FutureState,
\quad
Constraint,
\quad
Risk,
\quad
PredictiveEnvelope.
\]

BPCC 则利用自己的：

\[
SGC_{control}
+
FCS
+
z_{BPCC}
\]

在这些高层未来约束中完成高速局部控制。

因此可以用一句话概括整个设计：

\begin{statementbox}
身体负责持续拥有现实，高层认知被有意限制只能看清现实的一部分，
从而迫使其理解和预测整个世界；
这种高层理解再以未来结构约束的形式传给 BPCC，
由 BPCC 在完整局部现实中完成快速身体控制。
\end{statementbox}

进一步浓缩为 AgentOS 的一条核心多尺度原则：

\[
\boxed{
\begin{gathered}
\textbf{Upward:}\quad Detail\rightarrow Structure,\\
\textbf{Downward:}\quad Structure\rightarrow Constraint
\end{gathered}
}
\]

其中 SGC 认识孔径负责把现实细节逐步压缩成高层可理解结构，而 Kernel 的 Predictive Envelope 则把高层结构重新投影成 BPCC 可以实现的局部未来约束。二者共同构成从现实到认知、再从认知回到现实的双向结构梯度，也使 BPCC 不再只是一个快速控制器，而成为 AgentOS 多尺度认知—身体系统中承接高层预测并在现实中将其连续实现的身体动力学闭环。
